\begin{section}{Countable and Uncountable Sets} \label{sec:CountableUncountableSets}

A set $X$ is said to be \textbf{countable} if $\N \hyperref[thm:BijectionEquivalenceRelation]{\sim} X$.
A set $X$ is said to be \textbf{uncountable} if it is neither finite nor countable.
A set $X$ is said to be \textbf{at most countable} if $X$ is either finite or countable.

\bigskip

\begin{theorem} \label{thm:InfiniteSetBijection}
    Let $X$ be an infinite set. If $X \hyperref[thm:BijectionEquivalenceRelation]{\sim} Y$, then $Y$ is infinite.
\end{theorem}

\begin{proof}
    Suppose $Y$ is finite. Then there exists a bijection $f: J_{n} \to Y$ for some $n \in \N$.
    Since $X \hyperref[thm:BijectionEquivalenceRelation]{\sim} Y$, there exists a bijection $g: X \to Y$.
    \\[6mm]
    Then the composition $g^{-1} \circ f: J_{n} \to X$ is a bijection, which implies that $X$ is finite.
    This contradicts the assumption that $X$ is infinite. Thus, $Y$ is infinite.
\end{proof}

\bigskip

\begin{theorem} \label{thm:CountableSetBijection}
    Let $X$ be a countable set. If $X \hyperref[thm:BijectionEquivalenceRelation]{\sim} Y$, then $Y$ is countable.
\end{theorem}

\begin{proof}
    Since $X$ is countable, there exists a bijection $f: \N \to X$.
    Since $X \hyperref[thm:BijectionEquivalenceRelation]{\sim} Y$, there exists a bijection $g: X \to Y$.
    Then, by \Cref{thm:PropertiesOfCompositionFunctions}, the composition $g \circ f: \N \to Y$ is a bijection,
    which implies that $Y$ is countable.
\end{proof}

\bigskip

\begin{lemma} \label{lem:NIsInfinite}
    $\N$ is infinite.
\end{lemma}

\begin{proof}
    Suppose $\N$ is finite. Then there exists a bijection $f: J_{n} \to \N$ for some $n \in \N$.
    Then, by \Cref{thm:SubsetOfFiniteSetIsFinite}, $f(J_{n})$ is finite. By \Cref{thm:FiniteOrderedSetHasLeastAndGreatestElement},
    $f(J_{n})$ has a greatest element, say $m$.
    \\[6mm]
    Then we have $m + 1 \in \N$ and $m < m + 1 \notin f(J_{n})$,
    which contradicts the assumption that $f$ is surjective. Thus, $\N$ is infinite.
\end{proof}

\bigskip

\begin{theorem} \label{thm:CountableSetIsInfinite}
    Every countable set is infinite.
\end{theorem}

\begin{proof}
    Let $X$ be a countable set. Then there exists a bijection $f: \N \to X$. By
    \Cref{thm:InfiniteSetBijection,lem:NIsInfinite}, $X$ is infinite.
\end{proof}

\bigskip

\begin{remark}
    Let $X$ be a countable set. Then there exists a bijection $f: \N \to X$. The elements of $X$
    can be listed as $X = \set{x_1, x_2, x_3, \dots}$, where $x_n = f(n)$ for all $n \in \N$.
    In other words, the elements of $X$ can be arranged in a sequence.
\end{remark}

\bigskip

\begin{theorem} \label{thm:ZIsCountable}
    $\Z$ is countable.
\end{theorem}

\begin{proof}
    Let $f: \N \to \Z$ be defined by
    \[
    f(n) = \begin{dcases}
        -\frac{n}{2}, & \text{if } n \text{ is even}, \\
        \frac{n - 1}{2}, & \text{if } n \text{ is odd}.
    \end{dcases}
    \]
    We will show that $f$ is bijective, which implies that $\Z$ is countable.
    \begin{itemize}
        \item \textbf{Injectivity:} Suppose $f(n_1) = f(n_2)$ for some $n_1, n_2 \in \N$. Since $f(n) \geq 0$
        for odd $n$ and $f(n) < 0$ for even $n$, we have $n_1$ and $n_2$ are either both even or both odd.
            \begin{itemize}
                \item If both $n_1$ and $n_2$ are even, then $-\dfrac{n_1}{2} = -\dfrac{n_2}{2}$,
                which implies $n_1 = n_2$.
                \item If both $n_1$ and $n_2$ are odd, then $\dfrac{n_1 - 1}{2} = \dfrac{n_2 - 1}{2}$,
                which implies $n_1 = n_2$.
            \end{itemize}
            Thus, $f$ is injective.
        \item \textbf{Surjectivity:} Let $k \in \Z$. We need to find an $n \in \N$ such that $f(n) = k$.
            \begin{itemize}
                \item If $k \geq 0$, let $n = 2k + 1$. Then $f(n) = \dfrac{n - 1}{2} = \dfrac{2k}{2} = k$.
                \item If $k < 0$, let $n = -2k$. Then $f(n) = -\dfrac{n}{2} = -\dfrac{-2k}{2} = k$.
            \end{itemize}
            Thus, $f$ is surjective.
    \end{itemize}
    Since $f$ is both injective and surjective, we have $f$ is bijective.
\end{proof}

\bigskip

\begin{lemma} \label{lem:InfiniteSubsetOfNIsCountable}
    Every infinite subset of $\N$ is countable.
\end{lemma}

\begin{proof}
    Let $X \subseteq \N$ be an infinite subset. Let $X_{1} \coloneqq X$ and define a function $f: \N \to X$ as follows:
    \begin{itemize}
        \item $f(n) = \min X_{n}$
        \item $X_{n + 1} \coloneqq X_{n} \setminus \set{f(n)}$
    \end{itemize}
    for each $n \in \N$. By \Cref{thm:WellOrderingProperty}, $X_{n}$ has a least element, so $f$ is well-defined.
    We will show that $f$ is bijective, which implies that $X$ is countable.
    \begin{itemize}
        \item \textbf{Injectivity:} Suppose $n_{1} \neq n_{2}$ for some $n_{1}, n_{2} \in \N$. Without loss of generality,
        assume that $n_{1} < n_{2}$. Then $f(n_{1}) \in X_{n_{1}}$ and $f(n_{2}) \in X_{n_{2}}$. Since $X_{n_{2}}
        \subseteq X_{n_{1}} \setminus \set{f(n_{1})}$, we have $f(n_{2}) \neq f(n_{1})$. Thus, $f$ is injective.
        \item \textbf{Surjectivity:} Let $x \in X$. Let $S \coloneqq \setbuilder{n \in X}{n < x}$. If $S = \emptyset$,
        then $f(1) = x$. Otherwise, by \Cref{lem:SubsetOfJnIsFinite}, $S \subseteq \set{1, 2, \dots, x - 1}$ is finite.
        Let $k = |S|$, $S = \set{s_{1}, s_{2}, \dots, s_{k}}$ with $s_{1} < s_{2} < \dots < s_{k}$. Then we have
        \begin{itemize}
            \item $\min X_{1} = \min S \implies X_{2} = X \setminus \set{s_{1}}$.
            \item $\min X_{2} = \min S \setminus \set{s_{1}} \implies X_{3} = X \setminus \set{s_{1}, s_{2}}$.
            \item[] $\vdots$
            \item $\min X_{k} = \min S \setminus \set{s_{1}, s_{2}, \dots, s_{k - 1}} \implies X_{k + 1} = X \setminus S$.
        \end{itemize}
        Thus, $f(k + 1) = \min X_{k + 1} = x$. Therefore, $f$ is surjective.
    \end{itemize}
\end{proof}

\bigskip

\begin{theorem} \label{thm:InfiniteSubsetOfCountableSetIsCountable}
    Every infinite subset of a countable set is countable.
\end{theorem}

\begin{proof}
    Let $X$ be a countable set and let $Y \subseteq X$ be an infinite subset. Since $X$ is countable, there exists
    a bijection $f: \N \to X$. By \Cref{thm:PropertiesOfRestrictionFunctions}, $f|_{f^{-1}(Y)}^{Y}: f^{-1}(Y) \to Y$
    is bijective. Then, by \Cref{thm:InfiniteSetBijection}, $f^{-1}(Y) \subseteq \N$ is infinite.
    \\[6mm]
    By \Cref{lem:InfiniteSubsetOfNIsCountable}, $f^{-1}(Y)$ is countable. Then there exists a bijection $g: \N \to f^{-1}(Y)$.
    The composition $f|_{f^{-1}(Y)}^{Y} \circ g: \N \to Y$ is bijective, which implies that $Y$ is countable.
\end{proof}

\bigskip

\begin{theorem} \label{thm:SurjectionFromN}
    Let $X$ be infinite. If there exists a surjection $f: \N \to X$, then $X$ is countable.
\end{theorem}

\begin{proof}
    For each $k \in X$, let $S_{k} \coloneqq f^{-1}(\set{k})$. Since $f$ is surjective, $S_{k} \neq \emptyset$ and
    by \Cref{thm:WellOrderingProperty}, $S_{k}$ has a least element for all $k \in X$. Let $g: X \to \N$ be defined by
    \[
    g(k) = \min S_{k} \text{ for all } k \in X.
    \]
    Suppose $g(k_{1}) = g(k_{2})$ for some $k_{1}, k_{2} \in X \implies \min S_{k_{1}} = \min S_{k_{2}}$. Since
    $S_{i} \cap S_{j} = \emptyset$ for all $i, j \in X$ with $i \neq j$, we have $S_{k_{1}} = S_{k_{2}} \implies
    k_{1} = k_{2}$. Thus, $g$ is injective.
    \\[6mm]
    Since $g$ is injective, by \Cref{thm:PropertiesOfRestrictionFunctions}, $g|^{g(X)}: X \to g(X)$ is bijective.
    Since $X$ is infinite, $g(X) \subseteq \N$ is infinite by \Cref{thm:InfiniteSetBijection}. Then, by
    \Cref{lem:InfiniteSubsetOfNIsCountable}, $g(X)$ is countable, which implies that $X$ is countable.
\end{proof}

\bigskip

\begin{theorem} \label{thm:CountableUnionOfCountableSetsIsCountable}
    Let $\set{X_{n}}_{n \in \N}$ be a sequence of countable sets. Then $\displaystyle
    \bigcup_{n = 1}^{\infty} X_{n}$ is countable, i.e., any countable union of countable sets is countable.
\end{theorem}

\begin{proof}
    Since each $X_{n}$ is countable, there exists a bijection $f_{n}: \N \to X_{n}$ for each $n \in \N$.
    For each $n \in \N$, we define the followings:
    \begin{itemize}
        \item $t_{n} = \min\setbuilder*{k \in \N}{n \leq \dfrac{k(k + 1)}{2}}$.
        \item $s_{n} = n - \dfrac{t_{n}(t_{n} - 1)}{2}$.
    \end{itemize}
    \Cref{thm:BoundedSubsetOfZHasExtremum} guarantees that such $t_{n}$ exists.
    Let $g: \N \to \displaystyle \bigcup_{n = 1}^{\infty} X_{n}$ be defined by
    \[
    g(k) = f_{s_{k}}(t_{k} - s_{k} + 1) \text{ for all } k \in \N.
    \]
    Let $x \in \displaystyle \bigcup_{n = 1}^{\infty} X_{n}$. Then there exists an $m \in \N$ such that
    $x \in X_{m}$. Since $f_{m}$ is surjective, there exists a $l \in \N$ such that $f_{m}(l) = x$. Let
    \begin{itemize}
        \item $k = \dfrac{(l + m - 1)(l + m - 2)}{2} + m \in \N$.
        \item $T = \setbuilder*{z \in \N}{k \leq \dfrac{z(z + 1)}{2}}$.
    \end{itemize}
    We will show that $l + m - 2 \notin T$ and $l + m - 1 \in T$, which implies that
    \begin{itemize}
        \item $t_{k} = \min T = l + m - 1$.
        \item $s_{k} = k - \dfrac{t_{k}(t_{k} - 1)}{2} = m$.
        \item $g(k) = f_{s_{k}}(t_{k} - s_{k} + 1) = f_{m}(l + m - 1 - m + 1) = f_{m}(l) = x$.
    \end{itemize}
    Consider the following inequalities when $z = l + m - 2$ and $z = l + m - 1$:
    \begin{itemize}
        \item Let $z = l + m - 2$. Then we have
        \begin{align*}
            &k - \dfrac{(l + m - 2)(l + m - 1)}{2} \\
            = \; &\dfrac{(l + m - 1)(l + m - 2)}{2} + m - \dfrac{(l + m - 2)(l + m - 1)}{2} \\
            = \; &m > 0.
        \end{align*}
        \item Let $z = l + m - 1$. Then we have
        \begin{align*}
            &k - \dfrac{(l + m - 1)(l + m)}{2} \\
            = \; &\dfrac{(l + m - 1)(l + m - 2)}{2} + m - \dfrac{(l + m - 1)(l + m)}{2} \\
            = \; &m - \dfrac{(l + m - 1)(l + m) - (l + m - 1)(l + m - 2)}{2} \\
            = \; &m - (l + m - 1) = 1 - l \leq 0.
        \end{align*}
    \end{itemize}
    Thus $g$ is surjective, which implies that $\displaystyle \bigcup_{n = 1}^{\infty} X_{n}$
    is countable by \Cref{thm:SurjectionFromN}.
\end{proof}

\bigskip

\begin{corollary} \label{cor:FiniteUnionOfCountableSetsIsCountable}
    Let $X_{1}, X_{2}, \dots, X_{n}$ be countable sets. Then $\displaystyle \bigcup_{i = 1}^{n} X_{i}$ is countable,
    i.e., any finite union of countable sets is countable.
\end{corollary}

\begin{proof}
    Let $X_{n + 1}, X_{n + 2}, \dots$ be countable sets. Then, by \Cref{thm:CountableUnionOfCountableSetsIsCountable},
    $\displaystyle \bigcup_{n = 1}^{\infty} X_{n}$ is countable. Suppose that $\displaystyle \bigcup_{i = 1}^{n} X_{i}$
    is finite. Then $\displaystyle X_{1} \subseteq \bigcup_{i = 1}^{n} X_{i}$ is finite by
    \Cref{thm:SubsetOfFiniteSetIsFinite}, which contradicts \Cref{thm:CountableSetIsInfinite}. Thus,
    $\displaystyle \bigcup_{i = 1}^{n} X_{i}$ is infinite. Since $\displaystyle \bigcup_{i = 1}^{n} X_{i}
    \subseteq \bigcup_{n = 1}^{\infty} X_{n}$, by \Cref{thm:InfiniteSubsetOfCountableSetIsCountable},
    $\displaystyle \bigcup_{i = 1}^{n} X_{i}$ is countable. 
\end{proof}